\documentclass[11pt]{article}
\usepackage[T1]{fontenc}
\usepackage[utf8]{inputenc}
\usepackage{lmodern}
\usepackage[letterpaper,margin=0.78in]{geometry}
\usepackage{microtype}
\usepackage{booktabs,longtable,tabularx,array,multirow}
\usepackage{amsmath,amssymb}
\usepackage{xcolor}
\usepackage{enumitem}
\usepackage{titlesec}
\usepackage{fancyhdr}
\usepackage{lastpage}
\usepackage{hyperref}
\usepackage[most]{tcolorbox}
\usepackage{parskip}
\usepackage{needspace}
\usepackage{ragged2e}
\usepackage{setspace}

\definecolor{Navy}{HTML}{17365D}
\definecolor{Blue}{HTML}{2F75B5}
\definecolor{Teal}{HTML}{1F6D6A}
\definecolor{Gold}{HTML}{B38B2D}
\definecolor{Gray}{HTML}{5E6B75}
\definecolor{LightBlue}{HTML}{EEF4FA}
\definecolor{LightGold}{HTML}{FBF6E8}
\definecolor{LightGray}{HTML}{F4F5F7}
\definecolor{LightTeal}{HTML}{EDF7F6}

\hypersetup{colorlinks=true,linkcolor=Navy,urlcolor=Blue,citecolor=Teal,pdfauthor={OpenAI},pdfsubject={Value investing study and implementation materials}}
\setlist[itemize]{leftmargin=1.5em,itemsep=2pt,topsep=3pt}
\setlist[enumerate]{leftmargin=1.7em,itemsep=2pt,topsep=3pt}
\renewcommand{\arraystretch}{1.16}
\setlength{\tabcolsep}{5pt}
\titleformat{\section}{\Large\bfseries\color{Navy}}{\thesection}{0.6em}{}[\vspace{-0.25em}\color{Blue}\rule{\linewidth}{0.5pt}]
\titleformat{\subsection}{\large\bfseries\color{Teal}}{\thesubsection}{0.5em}{}
\titleformat{\subsubsection}{\normalsize\bfseries\color{Gray}}{\thesubsubsection}{0.5em}{}
\pagestyle{fancy}
\fancyhf{}
\fancyhead[L]{\small\color{Gray}\DocShortTitle}
\fancyhead[R]{\small\color{Gray}Value Investing Reference Set}
\fancyfoot[C]{\small\color{Gray}Page \thepage\ of \pageref*{LastPage}}
\renewcommand{\headrulewidth}{0.4pt}
\renewcommand{\footrulewidth}{0pt}

\newtcolorbox{principlebox}[1][]{enhanced,breakable,colback=LightBlue,colframe=Blue,boxrule=0.6pt,arc=2mm,left=2mm,right=2mm,top=1.5mm,bottom=1.5mm,title={#1},fonttitle=\bfseries}
\newtcolorbox{cautionbox}[1][]{enhanced,breakable,colback=LightGold,colframe=Gold,boxrule=0.6pt,arc=2mm,left=2mm,right=2mm,top=1.5mm,bottom=1.5mm,title={#1},fonttitle=\bfseries}
\newtcolorbox{modernbox}[1][]{enhanced,breakable,colback=LightTeal,colframe=Teal,boxrule=0.6pt,arc=2mm,left=2mm,right=2mm,top=1.5mm,bottom=1.5mm,title={#1},fonttitle=\bfseries}
\newtcolorbox{worksheetbox}[1][]{enhanced,breakable,colback=white,colframe=Gray,boxrule=0.45pt,arc=1mm,left=2mm,right=2mm,top=1.5mm,bottom=1.5mm,title={#1},fonttitle=\bfseries}
\newcommand{\checkbox}{\(\square\)}
\newcommand{\blankline}{\rule{0.96\linewidth}{0.35pt}}
\newcommand{\sourceRef}[1]{\textit{Source anchor: #1}}
\newcolumntype{Y}{>{\RaggedRight\arraybackslash}X}
\newcolumntype{C}[1]{>{\centering\arraybackslash}p{#1}}

\newcommand{\makedoccover}[4]{%
\begin{titlepage}
\thispagestyle{empty}
\vspace*{1.0in}
{\color{Navy}\rule{\linewidth}{2.2pt}}\\[0.55in]
{\Huge\bfseries\color{Navy} #1\par}
\vspace{0.18in}
{\Large\color{Teal} #2\par}
\vspace{0.5in}
\begin{tcolorbox}[colback=LightGray,colframe=Gray,boxrule=0.5pt,arc=2mm]
\textbf{Purpose.} #3
\end{tcolorbox}
\vfill
{\large Prepared from the revised edition of \textit{The Intelligent Investor}, with Jason Zweig's commentary.\par}
\vspace{0.12in}
{\small #4\par}
\vspace{0.25in}
{\small\color{Gray}\parbox{0.90\textwidth}{\centering Educational material only. Historical thresholds and examples are preserved for study; they are not individualized investment, tax, or legal advice.}}\par
\vspace{0.35in}
{\color{Navy}\rule{\linewidth}{2.2pt}}
\end{titlepage}
}

\newcommand{\DocShortTitle}{Investor Playbooks}
\begin{document}
\makedoccover{Defensive and Enterprising Investor Playbooks}{Checklists, gates, review cadence, and failure controls}{Turns the source's two investor modes into operational procedures.}{Standalone LaTeX study document}
\tableofcontents\newpage

\begin{cautionbox}[Source and modernization boundary]
The underlying source is the fourth revised edition, with Benjamin Graham's text updated in 1971--1972 and later commentary by Jason Zweig. Historical yields, tax rules, company examples, market levels, and dollar thresholds are therefore period-specific. This document preserves the source's framework and labels later implementation ideas separately rather than silently treating historical figures as current facts.
\end{cautionbox}

\section{Investor-Mode Selection}
\begin{worksheetbox}[Choose the operating mode]
\checkbox\ I want minimal maintenance and can accept market-level results.\\
\checkbox\ I can follow a written allocation through major declines.\\
\checkbox\ I do not want to read filings and normalize earnings for individual companies.\\
\textbf{If mostly checked:} Defensive mode is the default.\\[4pt]
\checkbox\ I can commit recurring research time every month.\\
\checkbox\ I can read financial statements and footnotes.\\
\checkbox\ I will maintain a written thesis, valuation, and sell discipline.\\
\checkbox\ I can tolerate long periods in which my approach looks unfashionable.\\
\textbf{If mostly checked:} Enterprising mode may be operationally feasible.
\end{worksheetbox}

\section{Defensive Investor Playbook}
\subsection{Before funding}
\begin{itemize}
\item \checkbox\ Define emergency liquidity outside the investment portfolio.
\item \checkbox\ Set a stock/bond target inside a policy band; 50/50 is the source's central reference, not a universal prescription.
\item \checkbox\ Decide how often to rebalance and what drift threshold triggers action.
\item \checkbox\ Select low-cost diversified vehicles or a screened list of leading companies.
\item \checkbox\ Automate periodic contributions.
\end{itemize}

\subsection{For direct-stock selection}
\begin{itemize}
\item \checkbox\ Adequate enterprise size and business durability.
\item \checkbox\ Current ratio and leverage meet a conservative standard.
\item \checkbox\ No earnings deficit across the required history.
\item \checkbox\ Long dividend record where the criterion is relevant.
\item \checkbox\ Multi-year earnings growth is satisfactory.
\item \checkbox\ P/E is moderate on normalized earnings.
\item \checkbox\ P/B or another asset-value measure provides adequate support.
\item \checkbox\ Combined valuation does not depend on heroic growth assumptions.
\end{itemize}

\subsection{Defensive maintenance cadence}
\begin{center}
\begin{tabularx}{\linewidth}{@{}p{0.20\linewidth}Y Y@{}}
\toprule
\textbf{Cadence} & \textbf{Action} & \textbf{Do not do}\\
\midrule
Each contribution & Add cash according to target allocation. & Change strategy because of headlines.\\
Quarterly & Confirm allocation drift and cash needs. & Revalue every holding from daily price moves.\\
Annually & Review policy, fees, diversification, and major business deterioration. & Turn annual review into wholesale portfolio replacement.\\
After life change & Reassess liquidity, horizon, and allocation. & Use market level alone as the reason to redefine risk capacity.\\
\bottomrule
\end{tabularx}
\end{center}

\section{Enterprising Investor Playbook}
\subsection{Gate 1: Is there a defined edge?}
\begin{itemize}
\item \checkbox\ The opportunity arises from price/value dislocation, not merely recent price momentum.
\item \checkbox\ The analysis can be explained with elementary arithmetic and business facts.
\item \checkbox\ The expected advantage remains after conservative assumptions and transaction costs.
\item \checkbox\ The thesis identifies what the market may be mispricing and why that mispricing could persist temporarily.
\end{itemize}

\subsection{Gate 2: Financial quality}
\begin{itemize}
\item \checkbox\ Liquidity is adequate for the business model.
\item \checkbox\ Debt cannot force an equity holder into a bad outcome under a plausible downturn.
\item \checkbox\ Earnings are not dependent on one peak year or aggressive adjustments.
\item \checkbox\ Management's acquisitions, repurchases, and compensation are economically coherent.
\end{itemize}

\subsection{Gate 3: Valuation and margin of safety}
\begin{itemize}
\item \checkbox\ At least two valuation approaches point to a discount.
\item \checkbox\ A bear-case value is estimated.
\item \checkbox\ Purchase price leaves a meaningful cushion to the base-case value.
\item \checkbox\ The thesis can survive a delay in market recognition.
\item \checkbox\ Position size assumes the analysis may be wrong.
\end{itemize}

\section{Negative Checklist: Automatic Review Triggers}
Any checked item requires either rejection or a written exception memo.
\begin{itemize}
\item \checkbox\ New issue / IPO whose appeal depends on promotional growth claims.
\item \checkbox\ High-yield or lower-quality debt purchased mainly for coupon income.
\item \checkbox\ Company expanding primarily through debt-financed acquisitions.
\item \checkbox\ Complex capital structure with large convertible or warrant dilution.
\item \checkbox\ Headline EPS depends heavily on adjustments, exclusions, or one-time gains.
\item \checkbox\ Valuation assumes a high growth rate for more than a few years.
\item \checkbox\ Management repurchases stock near peak valuation while insiders sell heavily.
\item \checkbox\ Position requires leverage or margin to make the expected return attractive.
\item \checkbox\ Thesis is ``the price has already gone up, so it will keep going up.''
\end{itemize}

\section{Speculation Containment Protocol}
The source allows that a person may speculate, but insists that it be recognized and separated.
\begin{enumerate}
\item Establish a fixed speculative-capital ceiling before selecting the security.
\item Use a separate account or at minimum a distinct tracking bucket.
\item Never replenish the bucket merely because losses occurred.
\item Do not count speculative gains as evidence that the method is investment-grade.
\item Do not use margin to magnify the bucket.
\item Report speculative performance separately from the long-term investment policy.
\end{enumerate}

\section{Portfolio Review Checklist}
\begin{worksheetbox}[Annual review]
\textbf{Policy}\\
\checkbox\ Allocation still matches horizon, liquidity, and behavior.\\
\checkbox\ Rebalancing policy was followed.\\
\checkbox\ Speculation remained within its cap.\\[4pt]
\textbf{Costs}\\
\checkbox\ Fund expense ratios and advisory fees remain reasonable.\\
\checkbox\ Trading frequency did not rise without a documented reason.\\[4pt]
\textbf{Direct holdings}\\
\checkbox\ Balance sheets remain sound.\\
\checkbox\ Normalized earning power has not materially deteriorated.\\
\checkbox\ Management capital allocation remains owner-oriented.\\
\checkbox\ Original margin of safety has not been eliminated by price appreciation or business deterioration.\\[4pt]
\textbf{Behavior}\\
\checkbox\ No purchase was made because of fear of missing out.\\
\checkbox\ No sale was made solely because a sound holding declined in quoted price.\\
\checkbox\ Forecast confidence did not replace valuation discipline.
\end{worksheetbox}

\section{Sell / Hold Decision Tree}
\begin{enumerate}
\item \textbf{Has the thesis broken?} If yes, reassess value immediately; selling may be warranted.
\item \textbf{Has financial strength deteriorated?} If leverage or liquidity crosses the safety boundary, treat it as a new security.
\item \textbf{Has price reached or exceeded conservative value?} If yes, compare after-tax expected return with alternatives rather than selling mechanically.
\item \textbf{Is the only reason the quote fell?} If yes, return to business value and the original thesis.
\item \textbf{Is there a clearly superior opportunity?} Compare expected return, risk, tax cost, and analytical confidence.
\item \textbf{Has position size become excessive?} Rebalance concentration even if the business remains excellent.
\end{enumerate}

\section{Adviser Due-Diligence Checklist}
\begin{itemize}
\item \checkbox\ Compensation model and all-in fees are understandable.
\item \checkbox\ Performance claims are benchmarked, time-matched, and net of fees.
\item \checkbox\ Adviser can explain downside controls and behavior during past stress periods.
\item \checkbox\ Custody arrangements are independent and transparent.
\item \checkbox\ The investment process does not depend on opaque forecasts or constant trading.
\item \checkbox\ Recommendations are compatible with the investor's actual knowledge and willingness to monitor them.
\end{itemize}

\section*{Primary Source}
{\small Benjamin Graham, \textit{The Intelligent Investor}, revised edition with commentary by Jason Zweig and preface/appendix material by Warren E. Buffett.}

\end{document}
